\chapter{Alur Kerja Penjualan (Sales Workflow)}

Modul ini adalah jantung operasional perusahaan, memuat siklus lengkap mulai dari prospek hingga pengiriman barang. Siklus ini sangat terstruktur: \textit{Inquiry} $\rightarrow$ \textit{Quotation} $\rightarrow$ \textit{Sales Order} $\rightarrow$ \textit{Invoice} $\rightarrow$ \textit{Delivery Order}.

\section{Inquiry (Pemesanan/Pertanyaan)}
Permintaan awal sering kali datang dari \textit{website} eksternal dan terekam di sini.

\subsection{Langkah Menindaklanjuti Inquiry}
\begin{enumerate}
    \item Buka menu \textbf{Inquiries}.
    \item Pilih data pelanggan yang berstatus \textbf{Pending}.
    \item Hubungi pelanggan via WhatsApp. Jika sepakat, staf (Admin) menekan tombol \textbf{Buat Quotation}.
\end{enumerate}

\section{Quotation (Penawaran Harga)}
Sistem memfasilitasi pembuatan penawaran harga resmi dengan perhitungan diskon dan PPN (11\%).

\subsection{Membuat Quotation Baru}
\begin{enumerate}
    \item Masuk ke menu \textbf{Quotations} lalu klik \textbf{New Quotation}.
    \item Pilih \textbf{Customer} dari \textit{dropdown}.
    \item Pada tabel \textit{Items}, klik \textbf{Add item} dan pilih produk yang akan ditawarkan.
    \item Sesuaikan \textit{Quantity} dan berikan \textit{Diskon} (jika ada). Subtotal akan dihitung secara otomatis.
    \item Klik \textbf{Create}. Dokumen akan berstatus \textit{Draft}.
    \item Admin dapat mengklik tombol \textbf{Print PDF} untuk mengunduh penawaran resmi.
\end{enumerate}

\section{Sales Order \& Invoice}
Jika penawaran disetujui, Quotation dikonversi menjadi \textit{Sales Order} (SO). SO harus di-approve sebelum dapat diubah menjadi tagihan (\textit{Invoice}).

\subsection{Skenario Persetujuan Owner \& Tagihan}
\begin{enumerate}
    \item Admin meng-klik aksi \textbf{Konversi ke Sales Order} pada halaman \textit{Quotation}.
    \item Dokumen SO baru terbentuk dengan status \textit{Pending}.
    \item Pengguna dengan peran \textbf{Owner} membuka SO tersebut dan mengklik tombol \textbf{Approve}.
    \item Setelah disetujui, tombol \textbf{Buat Invoice} akan muncul.
    \item Admin mengklik tombol tersebut untuk menerbitkan tagihan.
    \item Tagihan secara otomatis memiliki \textit{ULID} enkripsi untuk tautan unduhan PDF \textit{Secure Link}.
\end{enumerate}

\section{Pembayaran Invoice (Kasir)}
\begin{enumerate}
    \item \textbf{Kasir} masuk ke menu \textbf{Invoices}.
    \item Kasir membuka Invoice berstatus \textit{Belum Bayar} atau \textit{Sebagian}.
    \item Kasir mengklik tombol aksi berwarna hijau \textbf{Catat Pembayaran}.
    \item Sistem akan menampilkan *Modal form*. Masukkan nominal transfer pada *field* \textbf{Jumlah}.
    \item \textit{Sistem Mencegah Kesalahan:} Anda tidak dapat menginput jumlah melebihi nominal di \textbf{Sisa Bayar} berkat validasi \textit{maxValue}.
    \item Setelah sukses, status berubah menjadi \textit{Lunas} jika sisa tagihan mencapai nol.
\end{enumerate}

\begin{figure}[H]
    \centering
    \includegraphics[width=0.9\textwidth]{placeholder_catat_pembayaran.png}
    \caption{Pop-up Form Catat Pembayaran dengan Validasi Maksimal}
    \label{figure:catat_pembayaran}
\end{figure}

\section{Delivery Order (Surat Jalan)}
Setelah barang disiapkan atau invoice dibayar (bergantung kebijakan perusahaan), bagian Gudang mengeluarkan Surat Jalan. Fitur ini diawasi ketat oleh peran \textbf{Gudang}.

\subsection{Eksekusi Surat Jalan}
\begin{enumerate}
    \item Bagian \textbf{Gudang} mengklik menu \textbf{Delivery Orders} $>$ \textbf{New Delivery Order}.
    \item Memilih \textbf{Sales Order} tujuan, \textbf{Gudang Asal}, dan \textbf{Supir}.
    \item Memasukkan nomor polisi kendaraan.
    \item Setelah di-Simpan, Surat Jalan PDF siap dicetak dan stok gudang otomatis terpotong.
\end{enumerate}
